\chapter{نگارش نخستین اسکریپت}

در فصل‌های گذشته، مجموعه‌ای غنی از ابزارهای خط فرمان را فراگرفتیم. اگرچه این ابزارها توانایی حل چالش‌های پردازشی متعددی را دارند، اما استفاده از آن‌ها کماکان به اجرای دستی و تک‌به‌تک در ترمینال محدود است. با ترکیب این ابزارها در قالب برنامه‌هایی که خود طراحی می‌کنیم، می‌توانیم «پوسته» (\en{Shell}) را برای اجرای خودکار توالی‌های پیچیده‌ای از وظایف توانمند سازیم. این فرآیند از طریق نگارش «اسکریپت‌های پوسته» (\en{Shell Scripts}) محقق می‌شود.

\section{ماهیت اسکریپت‌های پوسته}

به بیان ساده، اسکریپت پوسته فایلی متنی حاوی مجموعه‌ای از دستورات است. پوسته این فایل را خوانده و دستورات موجود در آن را به‌گونه‌ای اجرا می‌کند که گویی مستقیماً در خط فرمان وارد شده‌اند. پوسته ابزاری منحصربه‌فرد است، زیرا هم‌زمان هم یک رابط خط فرمان قدرتمند برای تعامل با سیستم‌عامل و هم یک مفسر (\en{Interpreter}) برای زبان‌های اسکریپت‌نویسی محسوب می‌شود. تقریباً تمامی عملیاتی که در خط فرمان قابل اجراست، در قالب اسکریپت نیز قابل پیاده‌سازی است و بالعکس. تاکنون تمرکز ما بر ویژگی‌هایی از پوسته بوده که مستقیماً در ترمینال کاربرد دارند؛ اما پوسته قابلیت‌های پیشرفته‌تری نیز دارد که به‌طور ویژه برای برنامه‌نویسی و اتوماسیون طراحی شده‌اند.

\section{فرآیند نگارش اسکریپت‌های پوسته}

برای ایجاد و اجرای موفق یک اسکریپت پوسته، پیمودن سه گام اساسی الزامی است:

\begin{enumerate}
  \item \textbf{تدوین اسکریپت:} اسکریپت‌های پوسته در قالب فایل‌های متنی ساده ذخیره می‌شوند؛ بنابراین برای نگارش آن‌ها به یک ویرایشگر متن (\en{Text Editor}) نیاز داریم. ویرایشگرهای استاندارد با ارائه قابلیت «برجسته‌سازی نحو» (\en{Syntax Highlighting})، اجزای مختلف کد را با رنگ‌های متمایز نمایش می‌دهند که این امر به شناسایی سریع خطاهای رایج کمک شایانی می‌کند. ابزارهایی نظیر \cmd{vim}، \cmd{gedit} و \cmd{kate} از گزینه‌های ایده‌آل برای این منظور هستند.

  \item \textbf{تخصیص مجوز اجرا:} سیستم‌عامل به‌طور پیش‌فرض اجازه اجرای فایل‌های متنی عادی را صادر نمی‌کند. برای تبدیل یک فایل متنی به یک اسکریپت عملیاتی، باید سطوح دسترسی (\en{Permissions}) آن را به‌گونه‌ای تنظیم کنیم که ویژگی «قابل اجرا بودن» (\en{Executable}) به آن اعطا شود.

  \item \textbf{استقرار در مسیرهای جست‌وجوی پوسته:} برای آنکه بتوانید اسکریپت‌های خود را بدون نیاز به ذکر مسیر کامل فراخوانی کنید، باید آن‌ها را در یکی از دایرکتوری‌های موجود در متغیر محیطی \cmd{PATH} قرار دهید. این متغیر شامل فهرستی از دایرکتوری‌هاست که پوسته به‌ترتیب در آن‌ها به جست‌وجوی فایل اجرایی می‌پردازد. برای استفاده‌ی شخصی، دایرکتوری \file{\textasciitilde/bin} گزینه‌ای مناسب است؛ در صورت عدم وجود، می‌توانید آن را ایجاد کرده و با افزودن \cmd{export PATH="\$HOME/bin:\$PATH"} به فایل \file{.bashrc}، مسیر آن را به \cmd{PATH} اضافه کنید. در مقابل، اگر اسکریپت قرار است برای تمام کاربران سیستم در دسترس باشد، مسیر استاندارد \file{/usr/local/bin/} انتخاب درستی است که به‌طور پیش‌فرض در \cmd{PATH} اکثر توزیع‌ها قرار دارد. با این کار، فراخوانی اسکریپت‌ها از هر نقطه‌ای در محیط پوسته به ساده‌ترین شکل ممکن میسر خواهد شد.
\end{enumerate}

\section{ساختار فایل اسکریپت}

طبق سنت دیرینه‌ی برنامه‌نویسی، برای نمایش ساختار یک اسکریپت بسیار ساده، برنامه‌ی «سلام دنیا» را ایجاد می‌کنیم. ویرایشگر متن خود را باز کرده و قطعه کد زیر را در آن وارد نمایید:

\begin{latin}
\begin{lstlisting}
#!/bin/bash

# This is our first script.

echo 'Hello World!'
\end{lstlisting}
\end{latin}

خط پایانی این اسکریپت کاملاً آشناست؛ این خط صرفاً دستور \cmd{echo} را به همراه یک آرگومان رشته‌ای فراخوانی می‌کند. خط دوم نیز ساختار آشنایی دارد؛ این خط یک «توضیح» (\en{Comment}) است، مشابه آنچه در فایل‌های پیکربندی سیستم مشاهده و ویرایش کرده‌ایم. نکته‌ی حائز اهمیت درباره‌ی توضیحات در اسکریپت‌های پوسته این است که آن‌ها می‌توانند در انتهای خطوط دستور نیز درج شوند، به شرط آنکه حداقل یک فاصله پیش از آن قرار گیرد؛ مانند نمونه‌ی زیر:

\begin{latin}
\begin{lstlisting}
echo 'Hello World!' # This is a comment too
\end{lstlisting}
\end{latin}

در فرآیند پردازش اسکریپت، هر عبارتی که پس از نماد \cmd{\#} قرار گیرد، توسط مفسر پوسته نادیده گرفته می‌شود.

همان‌طور که در بسیاری از موارد صادق است، این ساختار در خط فرمان نیز به درستی عمل می‌کند:

\begin{latin}
\begin{lstlisting}
[me@linuxbox ~]$ echo 'Hello World!' # This is a comment too
Hello World!
\end{lstlisting}
\end{latin}

اگرچه استفاده از توضیحات (\en{Comments}) در خط فرمان چندان رایج نیست، اما سیستم به خوبی آن‌ها را پشتیبانی می‌کند.

اما خط اول اسکریپت ما کمی رمزآلود به نظر می‌رسد. با وجود اینکه با نماد \cmd{\#} آغاز شده و در ظاهر شبیه به یک توضیح معمولی است، اما در واقع هدفی فراتر از آن را دنبال می‌کند. توالی نویسه‌های \cmd{\#!} ساختار ویژه‌ای به نام \en{shebang} (شِبَنگ) است. نقش \en{shebang} این است که به هسته‌ی سیستم‌عامل (\en{Kernel}) اعلام کند کدام مفسر (\en{Interpreter}) باید برای اجرای دستورات این اسکریپت فراخوانی شود. برای اجرای مستقیم اسکریپت به‌عنوان یک برنامه‌ی قابل‌اجرا (با استفاده از \cmd{./script.sh})، وجود این خط به عنوان نخستین سطر الزامی است. با این حال، اگر اسکریپت را به‌صورت \cmd{bash script.sh} اجرا کنید، نیازی به \en{Shebang} نخواهد بود؛ زیرا در این حالت، خودِ دستور \cmd{bash} مفسر را مشخص می‌کند.

بیایید فایل اسکریپت را با نام \file{hello\_world} ذخیره کنیم.

\section{تخصیص مجوزهای اجرایی}

گام بعدی در فرآیند آماده‌سازی، اعطای قابلیت اجرا به اسکریپت است. این عملیات به‌سادگی و با بهره‌گیری از دستور \cmd{chmod} محقق می‌شود:

\begin{latin}
\begin{lstlisting}
[me@linuxbox ~]$ ls -l hello_world
-rw-r--r-- 1 me me 63 2009-03-07 10:10 hello_world
[me@linuxbox ~]$ chmod 755 hello_world
[me@linuxbox ~]$ ls -l hello_world
-rwxr-xr-x 1 me me 63 2009-03-07 10:10 hello_world
\end{lstlisting}
\end{latin}

به‌طور معمول، دو سطح دسترسی استاندارد برای اسکریپت‌های پوسته تعریف می‌شود:

\begin{itemize}
  \item \cmd{755}: برای اسکریپت‌هایی که تمامی کاربران سیستم مجاز به اجرای آن‌ها هستند.
  \item \cmd{700}: برای اسکریپت‌هایی که منحصراً توسط مالک فایل قابل اجرا می‌باشند.
\end{itemize}

لازم به ذکر است که اسکریپت‌ها برای اجرا شدن حتماً باید دارای مجوز «خواندن» (\en{Read}) باشند؛ چرا که مفسر پوسته پیش از اجرای دستورات، نیازمند خواندن محتوای متنی فایل است.

\section{تعیین مسیر استقرار اسکریپت}

پس از تنظیم سطوح دسترسی، اکنون اسکریپت آماده‌ی بهره‌برداری است:

\begin{latin}
\begin{lstlisting}
[me@linuxbox ~]$ ./hello_world
Hello World!
\end{lstlisting}
\end{latin}

برای اجرای موفق، درج مسیر دقیق فایل پیش از نام آن الزامی است. در غیر این صورت، سیستم با خطای زیر مواجه خواهد شد:

\begin{latin}
\begin{lstlisting}
[me@linuxbox ~]$ hello_world
bash: hello_world: command not found
\end{lstlisting}
\end{latin}

علت این رخداد چیست؟ آیا تفاوتی میان اسکریپت ما و سایر ابزارهای سیستمی وجود دارد؟ در حقیقت خیر؛ اسکریپت فاقد هرگونه نقص فنی است، اما «مکان» استقرار آن برای سیستم ناشناخته است. همان‌طور که در مبحث متغیرهای محیطی اشاره شد، متغیر \cmd{PATH} نقش کلیدی در فرآیند مکان‌یابی فایل‌های اجرایی ایفا می‌کند.

سیستم‌عامل برای یافتن هر دستور اجرایی (در صورتی که مسیر مطلق آن ذکر نشده باشد)، فهرستی از دایرکتوری‌های پیش‌فرض را پیمایش می‌کند. به همین جهت، هنگام وارد کردن دستور \cmd{ls}، سیستم به‌طور خودکار مسیر \file{/bin/ls} را یافته و اجرا می‌کند. متغیر \cmd{PATH} در برگیرنده‌ی این فهرست از مسیرهاست که با نویسه‌ی دو‌نقطه (\cmd{:}) از یکدیگر تفکیک شده‌اند. با دستور زیر می‌توان محتویات جاری این متغیر را بازبینی کرد:

\begin{latin}
\begin{lstlisting}
[me@linuxbox ~]$ echo $PATH
/home/me/bin:/usr/local/sbin:/usr/local/bin:/usr/sbin:/usr/bin:/sbin:/bin:/usr/games
\end{lstlisting}
\end{latin}

چنانچه اسکریپت ما در یکی از این مسیرها مستقر باشد، بدون نیاز به ذکر آدرس دقیق اجرا خواهد شد. در فهرست فوق، دایرکتوری نخست یعنی \file{/home/me/bin/} حائز اهمیت است؛ اکثر توزیع‌های لینوکس به‌صورت پیش‌فرض دایرکتوری \file{bin} در دایرکتوری خانگی کاربر (\en{Home Directory}) را به \cmd{PATH} می‌افزایند تا بستر اجرای برنامه‌های شخصی کاربر فراهم گردد. با ایجاد این دایرکتوری و انتقال اسکریپت به آن، برنامه‌ی ما همانند یک ابزار استاندارد سیستمی رفتار خواهد کرد:

\begin{latin}
\begin{lstlisting}
[me@linuxbox ~]$ mkdir bin
[me@linuxbox ~]$ mv hello_world bin
[me@linuxbox ~]$ hello_world
Hello World!
\end{lstlisting}
\end{latin}

بدین ترتیب، اسکریپت با موفقیت در سلسله‌مراتب فایل‌های اجرایی کاربر ادغام شد.

چنانچه دایرکتوری مذکور در متغیر \cmd{PATH} لحاظ نشده باشد، می‌توان به سادگی آن را افزود. برای این کار، کافی است عبارت زیر را به فایل پیکربندی \file{.bashrc} اضافه کنید:

\begin{latin}
\begin{lstlisting}
export PATH="$HOME/bin:$PATH"
\end{lstlisting}
\end{latin}

این تغییر در تمامی نشست‌های (\en{Sessions}) آتی ترمینال اعمال خواهد شد. جهت اجرایی شدن تغییرات در نشست فعلی، بایستی فایل \file{.bashrc} بازخوانی شود؛ این فرآیند اصطلاحاً «\en{Source} کردن» نامیده می‌شود:

\begin{latin}
\begin{lstlisting}
[me@linuxbox ~]$ . .bashrc
\end{lstlisting}
\end{latin}

در اینجا، دستور نقطه (\cmd{.}) معادل فرمان \cmd{source} عمل می‌کند. این دستور داخلی پوسته، محتویات فایل تعیین‌شده را خوانده و دستورات موجود در آن را به‌گونه‌ای اجرا می‌کند که گویی مستقیماً از طریق صفحه‌کلید وارد شده‌اند.

\begin{note}
\textbf{نکته:} توزیع اوبونتو و اکثر مشتقات دبیان، به‌صورت هوشمند دایرکتوری \file{\textasciitilde/bin} را به متغیر \cmd{PATH} اضافه می‌کنند، به شرط آنکه این دایرکتوری در زمان اجرای فایل \file{.bashrc} (هنگام ورود به سیستم) وجود داشته باشد. بنابراین در این سیستم‌ها، کافی است دایرکتوری را ایجاد کرده و یک بار از حساب کاربری خود خارج و مجدداً وارد شوید (\en{Log out/in}).
\end{note}

\section{دایرکتوری‌های استاندارد استقرار اسکریپت‌ها}

مسیر \file{\textasciitilde/bin} گزینه‌ای ایده‌آل برای نگهداری اسکریپت‌هایی است که منحصراً برای کاربری شخصی طراحی شده‌اند. چنانچه اسکریپتی تدوین کنید که دسترسی به آن برای تمامی کاربران سیستم الزامی باشد، طبق سنت سیستم‌های یونیکسی، باید آن را در مسیر \file{/usr/local/bin/} مستقر نمایید. همچنین، اسکریپت‌هایی که با هدف مدیریت سیستم و جهت استفاده کاربر ریشه (\en{Root}) طراحی می‌شوند، غالباً در دایرکتوری \file{/usr/local/sbin/} قرار می‌گیرند.

به‌طور کلی، نرم‌افزارهایی که به‌صورت محلی (\en{Locally}) توسعه یافته یا نصب می‌شوند — اعم از اسکریپت‌های پوسته یا برنامه‌های کامپایل‌شده — باید در زیرشاخه‌های \file{/usr/local/} سازمان‌دهی شوند و نه در دایرکتوری‌های \file{/bin} یا \file{/usr/bin}. مطابق با «استاندارد سلسله‌مراتب سیستم فایل لینوکس» (\en{FHS})، این دو دایرکتوری صرفاً به فایل‌هایی اختصاص دارند که توسط توزیع‌کننده (\en{Vendor}) تأمین و مدیریت می‌شوند.

\section{بهینه‌سازی ساختار و قالب‌بندی کد}

یکی از اهداف بنیادین در اسکریپت‌نویسی حرفه‌ای، تضمین «قابلیت نگهداری» (\en{Maintainability}) است؛ به این معنا که نویسنده یا سایر توسعه‌دهندگان بتوانند به سهولت اسکریپت را جهت انطباق با نیازمندی‌های جدید، بازبینی و اصلاح کنند. ارتقای خوانایی و شفافیت کد، از حیاتی‌ترین راهکارها برای تسهیل فرآیند نگهداری و توسعه‌ی پایدار اسکریپت محسوب می‌شود.

\subsection{گزینه‌های بلند در فراخوانی آرگومان‌ها}

بسیاری از دستوراتی که تاکنون بررسی کرده‌ایم، از هر دو قالب گزینه‌های کوتاه (\en{Short Options}) و گزینه‌های بلند (\en{Long Options}) پشتیبانی می‌کنند. برای نمونه، دستور \cmd{ls} دارای گزینه‌های متعددی است که می‌توان آن‌ها را به هر دو شکل فراخوانی کرد؛ مثال زیر را در نظر بگیرید:

\begin{latin}
\begin{lstlisting}
[me@linuxbox ~]$ ls -ad
\end{lstlisting}
\end{latin}

این دستور دقیقاً با عبارت زیر معادل است:

\begin{latin}
\begin{lstlisting}
[me@linuxbox ~]$ ls --all --directory
\end{lstlisting}
\end{latin}

در محیط ترمینال و تعامل مستقیم با خط فرمان، استفاده از گزینه‌های کوتاه جهت افزایش سرعت و کاهش حجم تایپ ترجیح داده می‌شود؛ اما در زمان نگارش اسکریپت‌های پوسته، استفاده از نام‌های طولانی گزینه‌ها به شکل قابل‌توجهی موجب ارتقای خوانایی و درک بهتر منطق برنامه می‌گردد.

\subsection{فاصله‌گذاری و شکستن خطوط دستور}

در مواجهه با دستورات طولانی، تقسیم فرمان به چندین سطر می‌تواند وضوح کد را به شکل چشم‌گیری افزایش دهد. در فصل ۱۷، نمونه‌ای نسبتاً پیچیده از دستور \cmd{find} را بررسی کردیم:

\begin{latin}
\begin{lstlisting}
[me@linuxbox ~]$ find playground \( -type f -not -perm 0600 -exec chmod 0600 ‘{}’ ‘;’ \) -or \( -type d -not -perm 0700 -exec chmod 0700 ‘{}’ ‘;’ \)
\end{lstlisting}
\end{latin}

بدیهی است که درک منطق این دستور در نگاه نخست دشوار است. در ساختار یک اسکریپت، می‌توان همین دستور را با فرمت زیر بازنویسی کرد تا خوانایی آن بهبود یابد:

\begin{latin}
\begin{lstlisting}
find playground \
\( \
    -type f \
    -not -perm 0600 \
    -exec chmod 0600 ‘{}’ ‘;’ \
\) \
-or \
\( \
    -type d \
    -not -perm 0700 \
    -exec chmod 0700 ‘{}’ ‘;’ \
\)
\end{lstlisting}
\end{latin}

با بهره‌گیری از تکنیک «ادامه‌ی خط» (ترکیب نویسه‌ی بک‌اسلش \cmd{\textbackslash} و خط جدید) و اِعمال «تورفتگی» (\en{Indentation})، ساختار سلسله‌مراتبی و منطق این دستور پیچیده برای بازبینِ کد شفاف‌تر می‌شود. هرچند این شیوه در خط فرمان نیز قابل اجراست، اما به دلیل دشواری در تایپ و ویرایش لحظه‌ای، به‌ندرت در محیط ترمینال استفاده می‌شود. شایان ذکر است که یکی از تفاوت‌های عملکردی میان محیط اسکریپت و خط فرمان، امکان استفاده از کاراکتر \en{Tab} جهت ایجاد تورفتگی در اسکریپت‌هاست؛ در حالی که در خط فرمان، این کلید منحصراً برای قابلیت «تکمیل خودکار» (\en{Auto-completion}) رزرو شده است.

\section{پیکربندی vim برای توسعه اسکریپت}

ویرایشگر متن \cmd{vim} از قابلیت‌های پیکربندی گسترده‌ای برخوردار است. تنظیم چند گزینه‌ی کلیدی می‌تواند فرآیند نگارش اسکریپت را به شکل مؤثری تسهیل نماید.

دستور زیر، قابلیت «برجسته‌سازی نحو» (\en{Syntax Highlighting}) را فعال می‌کند:

\begin{latin}
\begin{lstlisting}
:syntax on
\end{lstlisting}
\end{latin}

با اِعمال این تنظیم، المان‌های مختلف دستورات با رنگ‌های متمایز نمایش داده می‌شوند که علاوه بر جذابیت بصری، به شناسایی سریع خطاهای نگارشی کمک می‌کند. توجه داشته باشید که این ویژگی مستلزم نصب نسخه‌ی کامل \cmd{vim} است و فایل در حال ویرایش نیز باید دارای خط \en{shebang} باشد تا ماهیت اسکریپتی آن توسط ویرایشگر تشخیص داده شود. در صورت بروز مشکل، می‌توانید از دستور جایگزین زیر استفاده کنید:

\begin{latin}
\begin{lstlisting}
:set syntax=sh
\end{lstlisting}
\end{latin}

دستور زیر، ویژگی برجسته‌سازی نتایج جستجو را فعال می‌نماید:

\begin{latin}
\begin{lstlisting}
:set hlsearch
\end{lstlisting}
\end{latin}

به‌عنوان مثال، چنانچه واژه‌ی \cmd{echo} را جستجو کنید، با فعال بودن این گزینه، تمامی تکرارهای این واژه در متن متمایز خواهند شد.

دستور زیر، تعداد ستون‌های اختصاص یافته به کاراکتر \en{Tab} را تنظیم می‌کند:

\begin{latin}
\begin{lstlisting}
:set tabstop=4
\end{lstlisting}
\end{latin}

مقدار پیش‌فرض معمولاً ۸ ستون است. تنظیم این مقدار بر روی عدد ۴ (که یک رویه‌ی رایج در برنامه‌نویسی است)، موجب می‌شود تا خطوط طولانی با وضوح بیشتری در صفحه نمایش گنجانده شوند.

دستور زیر قابلیت «تورفتگی خودکار» (\en{Auto-indent}) را فعال می‌کند:

\begin{latin}
\begin{lstlisting}
:set autoindent
\end{lstlisting}
\end{latin}

این ویژگی سبب می‌شود که \cmd{vim} هنگام رفتن به خط جدید، تورفتگی را دقیقاً مطابق با خط پیشین تنظیم کند. این قابلیت سرعت تایپ ساختارهای شرطی و حلقه‌ها را به مراتب افزایش می‌دهد. جهت بازگشت تورفتگی به عقب (کاهش سطح تورفتگی)، می‌توانید از ترکیب کلیدهای \cmd{Ctrl+d} استفاده نمایید.

برای تثبیت و دائمی کردن این تنظیمات، کافی است دستورات فوق را (بدون درج نویسه‌ی دونقطه در ابتدای آن‌ها) در فایل پیکربندی شخصی خود به آدرس \file{\textasciitilde/.vimrc} درج نمایید.

\section{جمع‌بندی}

در این بخش از مبحث اسکریپت‌نویسی، فرآیند تدوین اسکریپت‌ها و نحوه اجرای استاندارد آن‌ها را در سیستم‌عامل بررسی کردیم. همچنین آموختیم که چگونه با بهره‌گیری از تکنیک‌های قالب‌بندی (\en{Formatting})، خوانایی کد را ارتقا دهیم؛ امری که مستقیماً موجب تسهیل در فرآیند نگهداری و توسعه‌ی آتی اسکریپت‌ها می‌گردد.

در فصول آینده، بارها با اصل «سهولت در نگهداری» (\en{Maintainability}) روبه‌رو خواهیم شد که به عنوان سنگ‌بنای اسکریپت‌نویسی حرفه‌ای مطرح می‌شود.

\section{منابع مطالعه تکمیلی}

برای مشاهده برنامه‌ها و مثال‌های «\en{Hello World}» در زبان‌های برنامه‌نویسی گوناگون، به منبع زیر مراجعه کنید:

\begin{latin}
\url{http://en.wikipedia.org/wiki/Hello_world}
\end{latin}

مقاله ویکی‌پدیا زیر، سازوکار \en{shebang} را با جزئیات بیشتری شرح می‌دهد:

\begin{latin}
\url{http://en.wikipedia.org/wiki/Shebang_(Unix)}
\end{latin}